\section{The original row products are integrable}
\label{sec:row-integrability}

This chapter supplies the additional \(\mathrm{RowPairInt}\) input
retained in DI1.  Its analytic assertion is slightly more general:
every original smooth cuspidal function is bounded, using its one
archimedean growth exponent for all derivative words.  We prove the
needed decay argument on the actual arithmetic components.  We do not
use the unreconstructed passage to the global adelic height, nor
identify the source's cyclic differential operators with generators
of the full enveloping center.

The named published input in the decay argument is the finite rational
Siegel covering and its rational-root setup, precisely as stated in
Miller--Schmid, \emph{On the rapid decay of cuspidal automorphic forms},
equations (2.2)--(2.8), printed pages 6--8 of the author preprint
\texttt{arXiv:1106.2149v1}.  Their Theorem A also assumes full center
finiteness through the meaning of \emph{automorphic form}.
We do not apply that theorem here.  RPI7--RPI10 give the complete
derivative argument appropriate to the stronger, uniform growth
already present in the source.  FV supplies the exact algebraic
group, finite classes, arithmetic subgroups and finite quotient volume.

Throughout, \(F=D.\mathrm{host}.F\), \(E=D.\mathrm{host}.E\),
\(\sigma=D.\mathrm{sp.cm.sig}\), \(V:\mathrm{HSp}(D)\), and
\(J=V.\mathrm{gram}\).  Write
\[
 G_{\mathbb A}=\mathrm{GA}(D,V),\qquad
 \Gamma=\mathrm{GQ}(V),\qquad
 \mathbf G=\operatorname{Res}_{F/\mathbb Q}U(J),\qquad
 G_\infty=\mathbf G(\mathbb R).
 \tag{RPI.1}
\]
These are the original objects, identified by FV1--FV4 and FV10.
The infinite family \(\theta\) satisfies the original
\(\mathrm{IsInfFam}(\theta)\).  Finite and infinite factors commute.
For a real Lie-algebra vector \(X\), \(R_X f(g)\) means
\(\left.\frac{d}{dt}f(g\exp(tX))\right|_{t=0}\).
Products of differential operators have the original right-fold order.
Every Haar measure or fundamental domain introduced on a unipotent
subgroup is auxiliary; the row measure and domain at the end are
unchanged.

% BEGIN RPI1
\paragraph{RPI1: the exact componentwise uniform growth.}
\label{claim:RPI1}
Let \(\phi:G_{\mathbb A}\to\mathbb C\) satisfy the original
\(\mathrm{IsAutFunSm}(\theta,V,\phi)\).  Choose a compact open subgroup
\(K\subset G_f\) fixing \(\phi\) on the right.  For each \(s\in G_f\),
put
\[
 f_s(u)=\phi(su),\qquad
 \Gamma_s=\{\gamma_\infty:\gamma\in U(J)(F),\
                              \gamma_f\in sKs^{-1}\}.
 \tag{RPI.2}
\]
Then \(f_s\) is a smooth function on \(\Gamma_s\backslash G_\infty\).
There is one integer \(N\ge0\), before every differential word \(D\),
such that
\[
       |R_D f_s(u)|\le C_{D,s}\|u\|_{\rm lin}^{N}
                         \quad(u\in G_\infty).
 \tag{RPI.3}
\]
Here the faithful linear norm is at least one and submultiplicative;
it is fixed before \(D,s,u\).  A fixed real basis of the Lie algebra
may be used for all words.  No full-center condition is used.

\emph{Proof.}
FV10 shows that the finite slice of the original open level contains
a compact open \(K\), proves that \(\Gamma_s\) is arithmetic, and
checks its left action on \(f_s\).  SBD1 identifies the source
curve-product smoothness with manifold smoothness and identifies its
one-place \(R_{\tau,X}\) with the indicated matrix-group derivatives.
Choose a real basis consisting of vectors in individual infinite
factors.  Finite products and real linear combinations of these
operators are exactly finite linear combinations of source words.

Use the rational representation on the underlying rational space of
\(E^n\), its dual, and a trivial line.  In fixed Euclidean coordinates
let \(\|u\|_{\rm lin}\) be the operator norm of this direct sum.
It is faithful, submultiplicative and at least one.  The actual
CM-type coordinates give, with a fixed constant \(c_0\),
\[
 \mathrm{normInf}(\theta,V,u)
 =\sum_{\tau,i,j}\bigl(|u_{\tau,ij}|+
                                  |(u_\tau^{-1})_{ij}|\bigr)
 \le c_0\|u\|_{\rm lin}.
 \tag{RPI.4}
\]
Indeed entries are bounded by operator norms, the dual controls the
inverse matrices, and the fixed change of real coordinates contributes
only a constant.  This is the comparison proved in SBD5.

The \(\mathrm{growth}\) field of the original definition states
\(\exists N\,\forall l\,\forall x\,\exists C\,\forall u\in G_\infty\).
Take \(x=s\), use the exact derivative identification, and then
(RPI.4).  The same \(N\) works for every word.  Finite linear
combinations of words retain that exponent.  Negative bounding
constants may be enlarged to nonnegative ones.  Rank zero is the
one-point group and is immediate.  This proof has not interchanged
\(\forall l\) with \(\exists N\), and has not replaced the fixed
finite component by a uniform adelic-height bound.
% END RPI1

% BEGIN RPI2
\paragraph{RPI2: the literal isotropic-frame subgroup is a unipotent radical.}
\label{claim:RPI2}
Let \(C\in M_{n,r}(E)\) satisfy the source
\(\mathrm{IsIsoFrame}(V,C)\), and set \(U=\operatorname{im}C\).
The original subgroup \(\mathrm{NI}(V,C)\) is exactly
\(N_U(\mathbb A_F)\), where \(N_U\) is the unipotent radical of the
\(F\)-parabolic stabilizing \(U\).  Its source rational subgroup
\(\mathrm{NIQ}(V,C)\) is exactly \(N_U(F)\).

\emph{Proof.}
We recall the actual matrix argument of CTC11, including the point
that the source takes an algebraic subgroup closure.  Choose a
rational dual frame \(C'\) with \(C^*JC'=1\).  Replacing it by
\(C'-\frac12C(C'^*JC')\) makes it isotropic and retains that pairing.
Complete the hyperbolic pair to an orthogonal decomposition
\[
 E^n=U\oplus W\oplus U',\qquad
 J_{\rm W}=\begin{pmatrix}0&0&I_r\\0&J_W&0\\I_r&0&0\end{pmatrix}.
 \tag{RPI.5}
\]
The rational change of basis is fixed and invertible.
In these coordinates the three source conditions on \(m=g-1\) are:
its first block column is zero; its last block row is zero; and
it sends \(U^\perp\) into \(U\).  The third condition is tested in
the source on rational vectors, which suffice because a rational
basis of \(U^\perp\) spans its adelic scalar extension.
Unitarity now determines all remaining blocks:
\[
 n(A,B)=
 \begin{pmatrix}
 I_r&-A^*J_W&B-\frac12A^*J_WA\\
 0&I&A\\
 0&0&I_r
 \end{pmatrix},
       \qquad B+B^*=0.
 \tag{RPI.6}
\]
Here \(A:U'\to W\); its entries and the skew-Hermitian entries of
\(B\) form finite-dimensional \(F\)-vector spaces.  Conversely every
matrix (RPI.6) satisfies all three source conditions.
Multiplication is
\[
 n(A,B)n(A',B')=
 n\bigl(A+A',B+B'+\tfrac12(A'^*J_WA-A^*J_WA')\bigr).
 \tag{RPI.7}
\]
Thus the defining set already is a subgroup; the source's
\(\mathrm{Subgroup.closure}\) adds no elements.  Its coordinates
and their inverse are polynomial and give the exact restricted-product
and real topologies.  Rational matrices have rational \(A,B\), and
conversely, proving the assertion about \(\mathrm{NIQ}\).

For completeness identify the radical, not just its abstract
isomorphism type.  The \(F\)-cocharacter acting by \(t,1,t^{-1}\)
on (RPI.5) has \(P(\lambda)\) equal to the stabilizer of \(U\):
unitarity makes preservation of \(U\) also preserve \(U^\perp\),
and these are precisely the block conditions allowing the limit
of \(\lambda(t)g\lambda(t)^{-1}\) at zero.  Its contracted subgroup
has identity diagonal blocks and is exactly (RPI.6).
Milne's Proposition 17.60 identifies this contracted subgroup
with the unipotent radical.  The hypothesis is a smooth connected
reductive \(F\)-group, already proved for this very \(U(J)\) in
FV1--FV2; characteristic zero and the \(F\)-defined cocharacter
are explicit.  This also proves normality of \(N_U\) in the
stabilizer directly from its block description.  The cases \(W=0\)
and \(r=1\) require no modification.
% END RPI2

% BEGIN RPI3
\paragraph{RPI3: a finite-level product domain for the original \(N_I\).}
\label{claim:RPI3}
For \(N_U\) from RPI2 and any compact open
\(K_N\subset N_U(\mathbb A_{F,f})\), put
\[
 \Lambda_N=\{\gamma_\infty:\gamma\in N_U(F),\
                                      \gamma_f\in K_N\}.
 \tag{RPI.8}
\]
Then \(N_U(\mathbb A_{F,f})=N_U(F)K_N\);
\(\Lambda_N\) is an arithmetic lattice in \(N_U(F_\infty)\);
and there is a Borel fundamental domain \(\Omega_N\) for its real
left quotient with compact closure such that
\[
                 \Omega_N\times K_N
 \tag{RPI.9}
\]
is a fundamental domain for the original
\(N_U(F)\backslash N_U(\mathbb A_F)\).
Product Haar measure on (RPI.9) is finite and positive.

\emph{Proof.}
In (RPI.6), choose \(F\)-bases for the \(A\)-space and the
skew-Hermitian \(B\)-space.  They give polynomial coordinates
\(F^d\), \(F_\infty^d\), and \(\mathbb A_{F,f}^d\) on the three
groups.  The additive approximation in SBD6, coordinate by
coordinate, says that \(F^d\) is dense in the finite adelic
coordinate space.  Every open coset \(yK_N\) therefore meets
\(N_U(F)\), which proves the asserted factorization.

The preimage of a compact open finite subgroup under rational
matrices is commensurable with an integral matrix subgroup, by
the compact-open intersection argument in FV10.  Its real
projection is discrete, since the corresponding rational matrices
in a product of two compact sets are finite by FV4.
SBD10, applied to the same rational unipotent matrix group,
gives a compact real quotient.  Its proof uses a rational central
flag and a bounded exponential coordinate box, so it also applies
to the present \(F\)-group by restriction of scalars.

Here is the measurable choice needed for (RPI.9).  A discrete
subgroup acts freely and properly by left translations.
Choose finitely many local sections covering its compact quotient,
with relatively compact images.  Assign a quotient point to the
first of these sections containing it; the resulting Borel
partition gives a Borel section whose image \(\Omega_N\) has
compact closure.  This is the elementary selection argument of
FV6.  It is an exact domain, or an a.e. domain if boundary
representatives are retained instead.

For any adelic \(n\), use the finite factorization to left
translate it rationally until its finite component lies in \(K_N\);
then use \(\Lambda_N\) to put its real component in \(\Omega_N\).
Conversely, if a rational left translate identifies two points
of (RPI.9), its finite component belongs to \(K_N\); its real
component is therefore in \(\Lambda_N\), and the defining
property of \(\Omega_N\) gives uniqueness.
The shear law (RPI.7) has Jacobian one for additive product
measure.  It proves unimodularity and the product Haar
description directly.  Compactness gives finiteness, and a
nonempty local section and the nonempty compact open \(K_N\)
give positivity.  Every continuous rationally invariant
integrand on this quotient is absolutely integrable, as also
proved in CTC11.  No convention assigning zero to a divergent
integral is involved.
% END RPI3

% BEGIN RPI4
\paragraph{RPI4: descent of the original \(N_I\) cusp integral.}
\label{claim:RPI4}
Suppose in addition to RPI1 that
\(\mathrm{IsCuspFun}(V,\phi)\) holds.  For every original isotropic
frame \(C\), every \(s\in G_f\), and every \(u\in G_\infty\),
\[
 \int_{(\Gamma_s\cap N_U(F_\infty))\backslash N_U(F_\infty)}
                  f_s(nu)\,dn=0.
 \tag{RPI.10}
\]
This follows from the original \(\mathrm{NI}\) condition; it
does not assume a \(\mathrm{ZCusp}\) condition.

\emph{Proof.}
Take
\[
 K_N=N_U(\mathbb A_{F,f})\cap sKs^{-1}.
 \tag{RPI.11}
\]
It is compact open in \(N_U(\mathbb A_{F,f})\).  Openness is
the subspace property.  Compactness follows because (RPI.6)
is a closed polynomial subgroup and \(sKs^{-1}\) is compact.
The lattice (RPI.8) is exactly \(\Gamma_s\cap N_U(F_\infty)\);
the rational embeddings are injective, so no additional
intersection or finite-index replacement has occurred.

Use the same \(\mathrm{NI}\) identified in RPI2, the product
domain from RPI3 and product Haar measure.  The universal
measure and domain quantifiers in \(\mathrm{IsCuspFun}\) permit
this choice.  At the argument \(su\) its integral is
\[
 \begin{split}
 0&=\int_{\Omega_N}\int_{K_N}
             \phi(n_\infty n_fsu)\,dn_f\,dn_\infty\\
  &=\operatorname{vol}(K_N)\int_{\Omega_N}
                                      f_s(n_\infty u)\,dn_\infty .
 \end{split}
 \tag{RPI.12}
\]
For the second equality, finite and infinite factors commute and
\(s^{-1}n_fs\in K\), so
\(n_\infty n_fsu=sn_\infty u(s^{-1}n_fs)\).
Right \(K\)-invariance gives precisely the displayed integrand.
Every integral is convergent: the integrand is continuous and
the domain has compact closure.  The finite positive volume
of \(K_N\) may be divided out.  The real quotient integral is
unchanged up to a positive normalization of Haar measure,
which proves (RPI.10) in every such normalization.
% END RPI4

% BEGIN RPI5
\paragraph{RPI5: every rational parabolic contains a suitable original \(N_I\).}
\label{claim:RPI5}
Let \(P\) be any proper \(\mathbb Q\)-parabolic of the actual
\(\mathbf G\), and \(N=R_u(P)\).  There is an original isotropic
frame \(C\) such that the group
\(\operatorname{Res}_{F/\mathbb Q}N_U\) in RPI2 is a nontrivial
normal subgroup of \(N\).  In particular it is an actual
subgroup of this \(N\), with its original rational points.

\emph{Proof.}
The checked cocharacter theorem used in SBD9 applies here:
Milne, \emph{Algebraic Groups} (2017), Theorem 25.1 and
Proposition 17.60.  FV1--FV4 and FV10 give the required
connected reductive rational group, and \(P\) is defined over
that same field.  Write \(P=P(\lambda)\), \(N=U(\lambda)\)
for a rational cocharacter.  Restriction-of-scalars adjunction
turns it into an \(F\)-cocharacter of \(U(J)\).
Let \(E^n=\bigoplus_m V_m\) be its integral weight decomposition.
Hermitian invariance pairs \(V_m\) only with \(V_{-m}\), and
the pairing is perfect.  Properness of \(P\) implies a
nonzero weight: if all weights were zero, \(\lambda\) would
act trivially on the faithful space and \(P\) would be the
whole group.  Thus its largest weight \(M\) is positive.

Put \(U=V_M\), \(U'=V_{-M}\), and
\(W=\bigoplus_{-M<m<M}V_m\).  These spaces are defined over
\(E\); choose an \(E\)-basis \(C\) for \(U\) and the paired
dual basis for \(U'\).  They give exactly (RPI.5), with
\(U,U'\) isotropic.  Thus \(C\) is an original
\(\mathrm{IsIsoFrame}\).

The condition defining \(P(\lambda)\) permits a block from
weight \(j\) to weight \(i\) only when \(i\ge j\).
It therefore preserves \(V_M=U\).  Its radical \(N\) has
identity diagonal blocks and strictly raises weights.
Every nonidentity block of \(n(A,B)\) in (RPI.6) also strictly
raises these weights: \(U'\to W\) has weight \(m+M>0\),
\(W\to U\) has weight \(M-m>0\), and \(U'\to U\)
has weight \(2M>0\).  Hence \(N_U\subset U(\lambda)=N\).
Moreover \(P(\lambda)\subset P_U\), the stabilizer of \(U\),
and \(N_U\) is normal in \(P_U\) by RPI2.  It is therefore
normal in \(N\).  Its skew-Hermitian central block has
positive dimension because \(r=\dim_EU>0\) and \(E/F\)
is quadratic, so it is nontrivial.  All identities are
polynomial over \(F\), and persist under restriction of
scalars and every real or adelic scalar extension.
% END RPI5

% BEGIN RPI6
\paragraph{RPI6: all real cusp integrals, and all their right derivatives.}
\label{claim:RPI6}
Under the hypotheses of RPI4, \(f_s\) is cuspidal for every
proper rational parabolic of \(\mathbf G\):
\[
       \int_{(\Gamma_s\cap N)\backslash N} f_s(ng)\,dn=0
                       \quad(g\in G_\infty).
 \tag{RPI.13}
\]
Here and below \(N\) denotes its real points when used in
an integral.  Every finite right derivative \(R_Df_s\)
has the same property and is still \(\Gamma_s\)-invariant.

\emph{Proof.}
SBD10 proves that the real quotient in (RPI.13) is compact
and has finite nonzero invariant measure.  Put
\(H=\operatorname{Res}_{F/\mathbb Q}N_U\subset N\) as in RPI5
and \(\Lambda=\Gamma_s\cap N\).  The subgroup \(H\) is rational,
normal and unipotent.  The intersections \(\Lambda\cap H\)
and the image \(\bar\Lambda\) in \(\bar N=N/H\) are lattices.
The first assertion is SBD10.  We give the detail for the
image that is sometimes hidden in a quotient argument.

In a rational faithful unipotent matrix realization, logarithm
and exponential are finite polynomials with fixed rational
coefficients.  Arithmetic matrices have bounded denominators:
an arithmetic subgroup is a finite union of rational translates
of a subgroup of an integral matrix group.  Consequently their
logarithms have bounded denominators in a rational basis of
\(\mathfrak n\).  Projection to
\(\mathfrak n/\mathfrak h\) preserves that property.
Its image lies in a fixed real lattice, and exponentiation
is a polynomial homeomorphism, so \(\bar\Lambda\) is discrete.
Choose a rational central flag for \(\mathfrak n\) compatible
with the ideal \(\mathfrak h\).  Such a flag is obtained by
successively choosing a central rational line in the remaining
ideal: the last nonzero term of
\([\mathfrak n,[\mathfrak n,\ldots,\mathfrak h]]\) lies in
\(\mathfrak h\cap Z(\mathfrak n)\).  Once that ideal is exhausted,
continue in the quotient.  This is the central-flag construction
in SBD10 with the indicated refinement.
For each rational basis vector, a sufficiently large integral
multiple has its exponential and inverse integral; finite index
puts a further positive power in \(\Lambda\).  Their projected
coordinates span the quotient.  The descending coordinate-box
argument of SBD10 therefore makes
\(\bar\Lambda\backslash\bar N\) compact.

These coordinates also justify the integration formula used
here.  Take a Borel section for the compact quotient
\(\bar\Lambda\backslash\bar N\), lift it locally to \(N\),
and integrate first along \(H\).  The fiber over a lift \(w\)
is \((\Lambda\cap H)\backslash H\), with integrand
\(h\mapsto f_s(hw g)\).  All these groups are unimodular.
Conjugation by \(N\) on \(\mathfrak h\) is unipotent and
has determinant one; hence the fiber Haar measure is
consistent when a section changes.  In the compatible central
coordinates, successive translations have triangular Jacobian
one, so ordinary product Fubini gives this quotient integration
formula.  Finite local sections of the compact quotients ensure
absolute convergence.  RPI4 says that every inner fiber integral
vanishes, for every \(wg\).  Thus (RPI.13) vanishes.

For the derivative assertion, choose a bounded real fundamental
domain in \(N\) as in SBD10 and the selection in RPI3.
On its compact closure times a compact neighborhood of \(g\),
every derivative of the smooth function \(f_s(ng)\) is
uniformly bounded.  Differentiation under the finite integral
is therefore justified by the mean-value estimate and dominated
convergence.  Differentiating the identically zero function
of \(g\) in (RPI.13), repeatedly, proves the assertion for
every ordered right word, with no reversal.  Right
differentiation commutes with left \(\Gamma_s\)-translation,
so the invariance assertion follows as well.
% END RPI6

% BEGIN RPI7
\paragraph{RPI7: a compact mean-zero derivative estimate.}
\label{claim:RPI7}
Let \(N_1\) be a real unipotent group and \(\Lambda_1\) an
arithmetic lattice in it.  Fix a real basis \(Y_1,\ldots,Y_a\)
of its Lie algebra.  There are a compact set
\(B_1\subset N_1\), containing one, and a constant \(c_1\)
such that, for every smooth \(\Lambda_1\)-invariant function
\(h\) whose quotient integral is zero,
\[
 |h(1)|\le c_1
       \max_{\substack{n\in B_1\\1\le j\le a}}
       \left|\left.\frac{d}{dt}h(\exp(tY_j)n)
                                                   \right|_{t=0}\right|.
 \tag{RPI.14}
\]
The compact set and constant depend only on the lattice and
basis, not on \(h\).  We will use this for positive-dimensional
radicals; the zero-dimensional case has \(h=0\).

\emph{Proof.}
Choose a Borel real fundamental domain \(\Omega_1\) with
compact closure.  Exponential and logarithm are inverse
polynomial maps by the argument in SBD10.  Set
\[
 B_1=\{\exp(t\log n):n\in\overline{\Omega_1},\ 0\le t\le1\}.
 \tag{RPI.15}
\]
It is compact and contains one.  Write
\(\log n=\sum_j a_j(n)Y_j\); the sum
\(\sum_j|a_j(n)|\) is bounded on \(\overline{\Omega_1}\),
say by \(c_1\).  The path \(t\mapsto\exp(t\log n)\)
has left velocity \(\log n\).  The fundamental theorem
of calculus therefore gives
\[
 |h(n)-h(1)|
 \le c_1\max_{q\in B_1,j}
       \left|\left.\frac{d}{dt}h(\exp(tY_j)q)\right|_{0}\right|.
 \tag{RPI.16}
\]
Normalize the quotient measure to mass one and integrate this
inequality over \(\Omega_1\).  Since \(\int h=0\),
\(h(1)=\int(h(1)-h(n))\,dn\), yielding (RPI.14).
This estimate does not presume that the displayed left fields
descend to the quotient: they are evaluated on the fixed compact
set of lifts \(B_1\), exactly where the path calculation takes
place.  Only left-lattice invariance of \(h\) and its genuine
compact quotient integral were used.
% END RPI7

% BEGIN RPI8
\paragraph{RPI8: the checked real reduction setup and its elementary bounds.}
\label{claim:RPI8}
For the actual group \(G_\infty\) and an arithmetic
\(\Gamma_s\), choose a rational minimal parabolic
\(P_0=MAN_0\) and positive simple rational roots \(\Delta\).
There are a compact \(C_\ell\subset MN_0\), a compact
\(C_r\subset G_\infty\), a number \(\epsilon>0\), and a
finite rational set \(Q\subset\mathbf G(\mathbb Q)\) with
\[
 \begin{gathered}
 \mathcal S=C_\ell A_\epsilon C_r,\qquad
 A_\epsilon=\{a\in A:\delta(a)\ge\epsilon
                                      \text{ for all }\delta\in\Delta\},\\
 G_\infty=\bigcup_{q\in Q}\Gamma_s q\mathcal S.
 \end{gathered}
 \tag{RPI.17}
\]
One may take \(C_r\) a maximal compact subgroup.  There
are constants \(c,d>0\), with \(d\) an integer, such that
\[
       \|a\|_{\rm lin}\le
         c\max(1,\max_{\delta\in\Delta}\delta(a))^d
                                      \quad(a\in A_\epsilon).
 \tag{RPI.18}
\]
For each \(\alpha\in\Delta\), there is a proper rational
parabolic whose radical \(N_\alpha\) has Lie algebra
\[
 \mathfrak n_\alpha=
   \bigoplus_{\substack{\beta>0\\
               \text{coefficient of }\alpha\text{ in }\beta>0}}
                                      \mathfrak g_\beta.
 \tag{RPI.19}
\]
It is normalized by \(MN_0\).  If \(r=\alpha(a)\ge1\), then
\[
 \|\operatorname{Ad}(a^{-1})|_{\mathfrak n_\alpha}\|
                                      \le c_\alpha r^{-1}.
 \tag{RPI.20}
\]

\emph{Published input and checks.}
We use precisely the rational-root and Siegel setup
(2.2)--(2.8) in Miller--Schmid, printed pages 6--8.
FV1--FV4 and FV10 identify the group as the real points of a
connected reductive rational matrix group; FV10 proves its
real center compact and proves arithmeticity of this exact
\(\Gamma_s\).  These are their standing hypotheses (2.1)--(2.2).
Their finite rational covering (2.8) gives (RPI.17).
Their chamber is written with all positive roots; passing
to a possibly smaller positive lower bound on the finitely
many simple roots, or enlarging the chamber, gives our form.
Enlarging a covering set preserves coverage.  No automorphic
form, center-finiteness or decay hypothesis enters this
published reduction statement.

\emph{Proof of the remaining assertions.}
Rational roots form a possibly nonreduced finite root system.
Every positive root is a nonnegative integral combination
of the simple roots.  These roots span the dual of
\(\operatorname{Lie}A\): a vector annihilated by them acts
trivially in the adjoint representation and would give a
split central direction, whereas the real center is compact.
Diagonalize the fixed faithful rational representation on
the split torus.  Each of its finitely many weights is a
rational linear combination of \(\Delta\).
For \(L_\delta=\log\delta(a)\), each \(L_\delta\ge\log\epsilon\).
If \(T=\max(0,\max_\delta L_\delta)\), the value of each
weight on \(\log a\) is at most \(dT+c'\), for some common
integer \(d\): positive coefficients are bounded using \(T\),
negative coefficients using the fixed lower bound.
Exponentiating and allowing the fixed diagonalizing basis
proves (RPI.18).

For (RPI.19), choose a rational vector in the cocharacter
space pairing positively with \(\alpha\) and zero with the
other simple roots; multiply it by a positive integer to
obtain a rational cocharacter.  Milne's contracted-subgroup
description, with the same reductive-group hypotheses as
RPI5, gives a proper parabolic and exactly (RPI.19).
Indeed its positive weights are precisely those roots having
positive \(\alpha\)-coefficient.  The minimal parabolic
normalizes this radical.  Explicitly \(M\) centralizes \(A\)
and preserves its root spaces; brackets with positive root
spaces retain positive \(\alpha\)-coefficient, so the
connected group \(N_0\) preserves their sum as well.

Write \(\beta=\sum_\delta b_\delta\delta\) in (RPI.19).
Then \(b_\alpha\ge1\) and, for \(r=\alpha(a)\ge1\),
\[
       \beta(a)^{-1}
       =r^{-b_\alpha}\prod_{\delta\ne\alpha}
                                          \delta(a)^{-b_\delta}
       \le \epsilon^{-\sum_{\delta\ne\alpha}b_\delta}r^{-1}.
 \tag{RPI.21}
\]
There are finitely many roots; in a fixed root-space basis
this proves (RPI.20), with equivalence of finite-dimensional
norms supplying any remaining constant.  If \(\Delta\) is
empty, \(A=\{1\}\) and \(\mathcal S\) is compact; this case
will be handled directly.
% END RPI8

% BEGIN RPI9
\paragraph{RPI9: a quantitative derivative iteration with compact factors.}
\label{claim:RPI9}
Let \(f\) be a smooth \(\Gamma_s\)-invariant function satisfying
all the cusp integrals in RPI6 and the uniform growth (RPI.3).
Fix \(\alpha\in\Delta\), and choose \(B_\alpha\) from RPI7 for
\((\Gamma_s\cap N_\alpha)\backslash N_\alpha\).  Put
\[
       C_j=B_\alpha^jC_\ell\quad(j\ge0).
 \tag{RPI.22}
\]
Each \(C_j\) is a fixed compact subset of \(MN_0\).
For every integer \(k\ge0\), there is a constant \(A_k\),
independent of \(a,l,r\), such that
\[
 \begin{split}
 |f(l a r)|
  &\le A_k\alpha(a)^{-k}
       \max_{\substack{D\text{ a basis word}\\ |D|=k}}
       \ \sup_{\substack{l'\in C_k\\r'\in C_r}}
                             |R_Df(l'a r')|\\
  &\le A'_k\alpha(a)^{-k}\|a\|_{\rm lin}^{N}
 \end{split}
 \tag{RPI.23}
\]
whenever \(l\in C_\ell\), \(r\in C_r\),
\(a\in A_\epsilon\), and \(\alpha(a)\ge1\).
The exponent \(N\) is the one fixed before all derivative
words in RPI1.

\emph{Proof.}
Take a fixed real basis \(X_1,\ldots,X_d\) of
\(\mathfrak g\) and a root-space basis \(Y_j\) of
\(\mathfrak n_\alpha\).  For \(l\in C_j\), the function
\(n\mapsto R_Df(nlar)\) has zero integral on the compact
\((\Gamma_s\cap N_\alpha)\)-quotient by RPI6.
Apply RPI7 to it.  At a point \(nlar\), \(n\in B_\alpha\),
the left derivative in direction \(Y_j\) is the right
derivative in direction
\[
        \operatorname{Ad}(r^{-1}a^{-1}l^{-1}n^{-1})Y_j .
 \tag{RPI.24}
\]
This follows by moving \(\exp(tY_j)\) to the right of
\(nlar\); it fixes the order and the inverse in (RPI.24).

The compact family \(l^{-1}n^{-1}\), for
\(l\in C_j,n\in B_\alpha\), normalizes
\(\mathfrak n_\alpha\), by RPI8.  Its action on that
finite-dimensional space has bounded coefficients.
Then (RPI.20) contributes a factor \(\alpha(a)^{-1}\).
Finally \(\operatorname{Ad}(r^{-1})\), \(r\in C_r\),
has uniformly bounded coefficients in the \(X_i\)-basis.
Consequently RPI7 gives, for a constant \(c_j\),
\[
 \sup_{\substack{l\in C_j\\r\in C_r}}
                |R_Df(lar)|
 \le c_j\alpha(a)^{-1}
       \max_i\sup_{\substack{l'\in C_{j+1}\\r'\in C_r}}
                                  |R_{X_i}R_Df(l'ar')|.
 \tag{RPI.25}
\]
The finite sum over the basis is absorbed into \(c_j\).
No derivative of a coefficient in (RPI.24) is taken: at
each stage (RPI.25) is applied anew to the smooth cusp
function \(R_Df\).  This is why enlarging \(C_j\) to
\(C_{j+1}\) is sufficient.

Iterate (RPI.25) \(k\) times, starting with the empty
word.  It prepends basis vectors to the word and yields
every required right-ordered word of length \(k\); the
constant is the product of the finitely many \(c_j\).
For the second line of (RPI.23), there are only finitely
many such words.  Their bounds in (RPI.3) have the same
exponent \(N\).  Submultiplicativity gives
\[
 \|l'a r'\|_{\rm lin}^{N}
 \le
 \left(\sup_{C_k}\|l'\|_{\rm lin}
       \sup_{C_r}\|r'\|_{\rm lin}\right)^N
                                            \|a\|_{\rm lin}^{N}.
 \tag{RPI.26}
\]
Both suprema are finite and may depend on \(k\), which is
harmless.  They cannot change \(N\).  Taking the maximum
of the finitely many derivative constants proves the
claimed estimate.
% END RPI9

% BEGIN RPI10
\paragraph{RPI10: rapid decay from the supplied uniform exponent.}
\label{claim:RPI10}
A smooth arithmetic-quotient function with the cuspidality
of RPI6 and uniform growth (RPI.3) decays faster than every
fixed negative power of \(\|\cdot\|_{\rm lin}\) on
\(\mathcal S\).  It is bounded on the whole real group.
This conclusion needs no finiteness under the enveloping
center, no \(K\)-finiteness, and no Hilbert realization.

\emph{Proof.}
If \(\Delta\ne\varnothing\), put
\(R(a)=\max(1,\max_{\delta\in\Delta}\delta(a))\).
The subset where \(R(a)=1\) is compact: each simple root
lies between its positive fixed lower bound and one, and
the simple roots form coordinates on the split real torus.
Thus \(C_\ell aC_r\) over that subset is compact and
smooth \(f\) is bounded there.

Otherwise choose \(\alpha\) attaining \(R(a)>1\).
There are only finitely many choices, so the constants in
RPI9 may be replaced by their maximum.  By (RPI.18),
for every \(k\),
\[
                |f(l a r)|
                  \le B_k R(a)^{-k+dN}.
 \tag{RPI.27}
\]
Also \(\|lar\|_{\rm lin}\le B R(a)^d\).
Given \(m\ge0\), choose an integer \(k\ge d(N+m)\).
It follows that
\[
        |f(g)|\le B'_m\|g\|_{\rm lin}^{-m}
                                \quad(g\in\mathcal S).
 \tag{RPI.28}
\]
On the compact part just separated, enlarge \(B'_m\)
by the maximum of \(|f(g)|\|g\|_{\rm lin}^{m}\).
If \(\Delta\) is empty the whole Siegel set is compact
and the same assertion follows immediately.

It remains to cover all real points; a bound on one
Siegel set alone would not suffice.  For each of the
finitely many rational \(q\) in (RPI.17), set
\(f_q(g)=f(qg)\).  It is invariant under
\(q^{-1}\Gamma_s q\), an arithmetic group.
Its right derivatives are \((R_Df)(qg)\); hence
(RPI.3) holds with the same \(N\), after multiplication
of the constants by \(\|q\|_{\rm lin}^{N}\).
For each rational parabolic radical \(N'\), conjugation
by \(q\) sends its integral to that of \(f\) over the
rational radical \(qN'q^{-1}\), with a positive Haar
factor.  All such integrals vanish by hypothesis.
Thus RPI7--RPI9 and the argument above apply to \(f_q\)
on the fixed \(\mathcal S\).  Only its lattice-dependent
constants change.  Each \(f_q\) is bounded there, and the
finite maximum of their bounds, together with left
\(\Gamma_s\)-invariance and (RPI.17), bounds \(f\) on
all of \(G_\infty\).

The independent argument is the compact mean-zero and
root-contraction argument discussed in Miller--Schmid
Section 5, printed page 28; (RPI.22)--(RPI.27) give its
iteration and the uniformity needed here explicitly.
It does not invoke their Theorem A under weaker hypotheses.
% END RPI10

% BEGIN RPI11
\paragraph{RPI11: boundedness of every original smooth cuspidal function.}
\label{claim:RPI11}
For every original \(D,V,\theta,\phi\),
\[
 \begin{gathered}
 \mathrm{IsInfFam}(\theta),\quad
 \mathrm{IsAutFunSm}(\theta,V,\phi),\quad
 \mathrm{IsCuspFun}(V,\phi)\\
 \Longrightarrow\quad
       \exists B\ge0\ \forall g\in G_{\mathbb A},\
                                      |\phi(g)|\le B.
 \end{gathered}
 \tag{RPI.29}
\]
The assertion concerns the original functions and their
given finite-level invariance; it supplies no uniform
bound over varying \(\phi\).

\emph{Proof.}
Choose a compact open \(K\) inside the original open
finite level, as in RPI1.  FV9 gives a finite set of
finite representatives \(s\), chosen after \(K\) and
before \(g\).  RPI1 gives smoothness and one exponent
for all derivatives of each \(f_s\).  RPI2--RPI6 prove
all its real rational-parabolic cusp integrals from
the source's actual \(\mathrm{NI}\) definition.
RPI10 therefore bounds each \(f_s\) on the whole real
group, including the rational translates in its real
Siegel covering.

For arbitrary \(g\), FV10 gives the exact factorization
\[
       g=\gamma s(\gamma_\infty^{-1}g_\infty)k,
       \qquad \gamma\in U(J)(F),\quad k\in K.
 \tag{RPI.30}
\]
The original rational left and finite right invariance
make \(\phi(g)=f_s(\gamma_\infty^{-1}g_\infty)\).
Take the maximum over the finite set of component bounds.
The zero function and the rank-zero group are included.
The field \(\mathrm{zfin}\) of the source definition has
not been used; uniform derivative growth is sufficient.
Neither the polynomial unitary reduction nor
\(\mathrm{growth\_iff}\) has entered this proof.
% END RPI11

% BEGIN RPI12
\paragraph{RPI12: products on every original finite-volume domain.}
\label{claim:RPI12}
Let \(\mu\) be any original Haar measure on \(G_{\mathbb A}\),
and \(\mathcal D\) any original fundamental domain for
\(\Gamma\) with respect to \(\mu\).  Let
\(\phi,\psi\) satisfy the hypotheses of RPI11.
For every fixed \(a,b\in G_{\mathbb A}\) and
\(z,w\in\mathbb C\), the function
\[
       g\longmapsto
            \overline{z\,\phi(ga)}\,w\,\psi(gb)
 \tag{RPI.31}
\]
is integrable on the unchanged \(\mathcal D\), both for
\(\mu\) and for every finite nonnegative scalar multiple
of \(\mu\).  Finite linear combinations of such functions
are integrable there as well.  More generally the same
conclusion holds whenever \(\phi,\psi\) are bounded
continuous functions on this original group.

\emph{Proof.}
FV8 proves \(\mu(\mathcal D)<\infty\) for every such
measure and a.e. domain, using precisely the original
rational subgroup.  RPI11 gives constants \(B_\phi,B_\psi\).
Right translation is continuous and ranges over the same
whole group, so (RPI.31) is continuous and has absolute
value at most
\(|z||w|B_\phi B_\psi\).
Its absolute integral is at most this constant times
\(\mu(\mathcal D)\), which is finite.
Multiplying the measure by a finite nonnegative scalar
preserves the bound, including scalar zero.  Finite
summation preserves integrability.

In particular the unmodified functions and every fixed
right translate lie in \(L^2(\mathcal D,\mu)\).
One could use Cauchy--Schwarz, but the displayed bounded
product estimate already proves exactly what is needed.
This proof does not identify or renormalize \(\mu\),
does not change the domain, and does not require a
Schur inner-product scalar to exist.
% END RPI12

% BEGIN RPI13
\paragraph{RPI13: the actual multiplicity spaces and finite normalized actions.}
\label{claim:RPI13}
Fix any original \(V,S,\theta,M:\mathrm{DSModel}(V,S)\),
level set \(J_0\) and level character \(\chi\).
For every
\(T\in\mathrm{Mult}(\theta,M,J_0,\chi)\)
and every original smooth tuple \(v\), the function \(T(v)\)
is bounded and continuous.  The same is true of every
value of the original good part, and of its actual
normalized finite residual translates.

\emph{Proof.}
The definition of \(\mathrm{Mult}\) is the algebraic
complex span of maps satisfying \(\mathrm{IsMultMap}\).
Thus \(T\) is a finite linear combination of such maps.
For each generator \(T_j\), its \(\mathrm{aut}\) and
\(\mathrm{cusp}\) fields say, literally for every \(v\),
\[
 \mathrm{IsAutFunSm}(\theta,V,T_j(v)),\qquad
 \mathrm{IsCuspFun}(V,T_j(v)).
 \tag{RPI.32}
\]
Apply RPI11 to these functions.  A finite linear combination
is bounded and continuous, with bound the sum of the
absolute coefficients times the individual bounds.
This argument uses the span definition directly; it does
not need an unproved equivalence asserting that all five
multiplicity-map clauses are already closed under spans.
The source good part is an intersection with
\(\mathrm{Mult}\), so the same argument applies to it.

For clarity retain the exact choice defining the finite
action.  For a residual element \(h\), if an original
\(\mathrm{IsLevelLift}\) exists, the chosen lift
\(\kappa_h\in G_f\) gives
\[
 A_h\varphi(x)=a_h^{-1}\varphi(x\kappa_h),\qquad
 a_h=\mathrm{levChi}(V,\beta,\mathrm{flagL}(L),\kappa_h).
 \tag{RPI.33}
\]
The source flag-character identity has \(|a_h|=1\).
Indeed \(\mathrm{flagL}\) replaces the residual subgroup
character by the trivial character, so its \(\mathrm{ext1}\)
factor is one.  Every remaining factor of the exact
\(\mathrm{levFn}\) is the image of a unit of a finite residue
field under one of the prescribed multiplicative characters.
If that unit has order \(m>0\), its complex image has
\(m\)-th power one and therefore modulus one.  Products
and inverses retain modulus one.  This is the elementary
argument in \texttt{NormLevChiFlagLEqOne}, independent of
row integrability; in particular \(a_h\ne0\).
If no lift exists, the literal \(\mathrm{actG0}\) definition
is zero.  Both branches give bounded continuous functions.
Even for any fixed nonunit scalar the boundedness argument
would apply, but (RPI.33) preserves the actual normalization.
The finite sum defining \(\mathrm{gaL}\) and the original
central idempotent \(e_0\) therefore also preserves these
bounded continuous values.  This proof leaves the
\(\mathrm{DSModel}\), its smooth-vector space, every
coefficient map and every finite lift unchanged.
% END RPI13

% BEGIN RPI14
\paragraph{RPI14: the complete original row-pair producer.}
\label{claim:RPI14}
Fix exactly the original data of DI1,
\[
          D,G,A_3,K_s,R,\mathscr F,c_G,X_0,Y,Z,C,s_{EE},
 \tag{RPI.34}
\]
where \(\mathscr F\) denotes the source \(\mathrm{Face3}\)
called \(F\) there, and retain the original
\(\mathrm{Pins}(R,\mathscr F,c_G)\) and
\(\mathrm{Holds}(R,\mathscr F)\).
For every \(b\in\{0,1\}\), every original parameter \(t\),
every \(w\) in the full original \(P_b\)-level
multiplicity space of \(R.MY(0)\), every \(c\) in its
original good part, every two original smooth tuples
\(v,v'\), and every \(g\in G\), the function
\[
 h\longmapsto
       \overline{A_{R.e_Gg}(c(v))(h)}\,w(v')(h)
 \tag{RPI.35}
\]
is integrable on exactly \(R.\mathrm{fd0}\) for exactly
\[
 \mu_b=
 \mathrm{ENNReal.ofReal}\!\left(
                 \mathrm{nP}(\mathscr F.\mathrm{lv2},b)^{-1}
                            \right)\,R.\nu_0.
 \tag{RPI.36}
\]
This is the entire \(\mathrm{RowPairInt}\) of DI.3a,
with no additional analytic premise.

\emph{Proof.}
The common record
\(\mathrm{Pins.char.base.models}\) supplies the original
\(\mathrm{IsInfFam}(R.\theta)\), and its \(\nu_0\) field
supplies both Haar measure and the original
\(\mathrm{GQ}(Y_c(D,0))\)-fundamental domain
\((R.\nu_0,R.\mathrm{fd0})\).
These are exactly the premises used in the pinned proof
of \texttt{rowProj\_holds}; none has been chosen anew.

The two level spaces in the statement are the literal
\(\mathrm{Mult}\) and good part appearing in
\texttt{Good\allowbreak Projector\allowbreak Smallest\allowbreak Space\allowbreak Forms.\allowbreak RowProj}.
More explicitly, put \(P_b=R.\mathrm{levR}(\mathscr F.\mathrm{lv2},b)\).
The level datum of \(w\) is
\(R.\mathrm{datYt}(t,0,P_b)\), while the good part containing
\(c\) uses \(R.\mathrm{dat0}(P_b)\); both use the same
\(R.\beta_Y(0)\), the corresponding \(\mathrm{levS}\) and
\(\mathrm{levChi}\), and the same model \(R.MY(0)\).
RPI13 therefore bounds \(c(v)\) and \(w(v')\).
There is no need to change \(t\), to replace \(R.MY(0)\)
by another realization, or to restrict to an irreducible
summand.  For the finite action in (RPI.35), use its
chosen lift and scalar in (RPI.33).  RPI12 proves the
result for \(R.\nu_0,R.\mathrm{fd0}\), including the
zero branch if no lift exists.

The scalar in (RPI.36) is finite.  In fact the original
\(\mathrm{nP}\) is positive, so it is a positive scalar
and is exactly the normalized measure in DI.3.
RPI12 applies to that scalar multiple without changing
the domain.  All choices of temporary compact levels,
finite class representatives and real domains were used
only to prove boundedness of the particular functions;
they have disappeared from (RPI.35)--(RPI.36).
Thus the result has the stated universal order in
\(b,t,w,c,v,v',g\).  No common threshold, new row basis,
or altered quotient measure is inserted.
% END RPI14

% BEGIN RPI15
\paragraph{RPI15: the first division output from its original premises.}
\label{claim:RPI15}
The conditional constructor DI1--DI21 now proves the first,
literal \(\mathrm{DivDef3}\) output of
\texttt{divTied3\_read} from its original stated premises.
Its additional \(\mathrm{RowPairInt}\) premise is discharged
by RPI14 on the same objects.  In particular the original
\(\mathrm{RowProj}\) metric identity is unconditional relative
to those original common data.

\emph{Proof.}
DI11 uses \(\mathrm{RowPairInt}\) exactly to integrate
the finitely many normalized action terms on
\((R.\mathrm{fd0},\mu_b)\).
RPI14 supplies all those terms, including the identity
and inverse residual elements.  Its unchanged measure
and domain allow the quotient right-translation equality
already proved in RQC4 and used in DI11.
The self-adjoint central coefficient identity
\(e(g^{-1})=\overline{e(g)}\) gives the finite projector
adjoint identity.  Since \(e_0c=c\), it gives
\[
 \int_{R.\mathrm{fd0}}\overline{c(v)}\,w(v')\,d\mu_b
 =
 \int_{R.\mathrm{fd0}}\overline{c(v)}\,
                      (e_0w)(v')\,d\mu_b.
 \tag{RPI.37}
\]
Both products are integrable by RPI12--RPI14.

DI12 treats the literal choice definition of \(\mathrm{ipM}\).
Equation (RPI.37) and integrability make the sets of
\(\mathrm{IsMultIP}\) scalar witnesses equal.  If empty,
both choices return zero.  If nonempty, the original
nonzero smooth tuples make the scalar unique.
Hence the exact row-projector identity follows, with no
existence assertion for a Schur scalar added to the hypotheses.
The remaining steps and terminal identities of DI10,
the same-object composite identifications DI13, and the
finite rational adjoint and theta-subrecord arguments
DI14--DI21 are unaffected.  Applying their sufficient
constructor gives the claimed first \(\mathrm{DivDef3}\).

The original five real-group predicates in DI.2, all
three readings, the same group tie, \(\mathrm{Pins}\),
\(\mathrm{Holds}\), rational row tests, positive operators,
geometry and stored weak-chain activation remain the
premises of that constructor.  The present proof does
not produce any of them.  It also does not prove the
stronger second output \(\mathrm{DivRows3}\) or its
actual theta square.  In particular it does not replace
the separate nonzero \(\mathrm{IsItW}\) chain by the
argument about a possible \(\mathrm{IsMultIP}\) scalar.
The universal source \(\mathrm{CuspRapidDecay}\) with its
global quotient-height conclusion is not asserted here.
The precise gain is the removal of the additional
row-integrability premise from the first-output consumer.
% END RPI15
